\subsection{Verfassen eines Sitzungsprotokolls}
\label{protokoll}
Um die Beschlüsse aus den Sitzungen des Fachschaftsrats festzuhalten, werden \textit{Sitzungsprotokolle} angefertigt. Dafür gibt es eine eigens erstellte \hyperref[latex-fsrprotocol]{\LaTeX-Dokumentenklasse}, die so gestaltet ist, dass auch Mitglieder ohne \LaTeX-Erfahrung damit arbeiten können.

Grundlagen zu \LaTeX{}, zum Kompilieren der PDF-Datei und zum Umgang mit Fehlermeldungen finden sich in \autoref{latex-uebersicht}, \autoref{latex-kompilieren} und \autoref{latex-fehler-finden}.

\subsubsection{Erstellen eines neuen Protokolls}
\label{protokoll-erstellen}
Auf dem Arbeitsrechner des Fachschaftsrats ist eine Vorlage für Sitzungsprotokolle hinterlegt. Mit ihr lässt sich ein neues Protokoll direkt im Dateimanager anlegen:

\notebox{Die folgende Anleitung gilt für den eingerichteten Arbeitsrechner des Fachschaftsrats. Auf eigenen Rechnern müssen zunächst die in \autoref{latex-eigener-rechner} beschriebenen Voraussetzungen erfüllt sein. Statt über das Kontextmenü wird dort eine leere \texttt{.tex}-Datei angelegt und mit den Codeblöcken aus \autoref{protokoll-schreiben} gefüllt.}

\begin{enumerate}
	% TODO: Nach Fertigstellung des Servers den Pfad zu Dokumente/Protokolle/JJJJ ändern.
	\item Im Dateimanager den Ordner \texttt{Dokumente/webspace/Protokolle/JJJJ} für das Jahr der Sitzung öffnen. Falls der Ordner noch nicht existiert, muss er zuerst angelegt werden.
	\item Mit einem Rechtsklick auf eine freie Fläche das Kontextmenü öffnen und
	      \textit{Neu erstellen} \textrightarrow{} \texttt{fsrprotocol.tex} auswählen.
	\item Die neue Datei nach dem Schema \texttt{JJJJ-MM-DD\_Sitzungsprotokoll\_Name.tex} umbenennen,
	      z.\,B. \texttt{2026-10-14\_Sitzungsprotokoll\_Max-Mustermann.tex}.
	\item Die Datei in einem Editor öffnen (auf dem Arbeitsrechner des Fachschaftsrats: VSCodium) und wie in \autoref{protokoll-schreiben} beschrieben bearbeiten.
\end{enumerate}

\subsubsection{Schreiben eines Protokolls}
\label{protokoll-schreiben}
Im ersten Teil des \LaTeX-Dokuments werden die Dokumentenklasse geladen und die Metadaten der Sitzung festgelegt.

\importantbox{Die Angaben zur Sitzung sind Pflicht: Fehlt eine davon, verweigert \LaTeX{} die \hyperref[latex-kompilieren]{Kompilierung}. Die dabei angezeigte Fehlermeldung ist in \autoref{latex-fehlermeldungen} beschrieben.}

\notebox{Text hinter einem \texttt{\%} ist ein Kommentar und erscheint nicht im PDF. Ein \texttt{\%} im Text wird als \mintinline{latex}{\%} geschrieben (siehe \autoref{latex-sonderzeichen}).}
\begin{Codeblock}[title={Anfang des Dokuments}, unbreakable]{latex}
	\documentclass{fsrprotocol}

	% --- Setze Informationen zu dem Dokument --- %
	% Hier stehen Informationen zur Sitzung. Der Protokollführer muss diese ausfüllen.
	\newcommand{\Sitzungsdatum}{dd.mm.yyyy}
	\newcommand{\Sitzungszeitraum}{00:00 - 23:59 Uhr}
	\newcommand{\Sitzungsort}{Hier Veranstaltungsort eintragen}
	\newcommand{\Sitzungsleitung}{Vorname Nachname}
	\newcommand{\Protokollfuehrer}{Vorname Nachname}
	\newcommand{\Anwesenheitsliste}{Vorname1 Nachname1, Vorname2 Nachname2, ...}
	\newcommand{\Gaesteliste}{Vorname1 Nachname1, Vorname2 Nachname2, ...}

\begin{document}
\end{Codeblock}

Nur die Platzhalter in den geschweiften Klammern der \mintinline{latex}{\newcommand}-Zeilen werden ersetzt; alles andere bleibt unverändert. Waren keine Gäste anwesend, wird bei \mintinline{latex}{\Gaesteliste} \enquote{keine} eingetragen. Hintergründe zu \mintinline{latex}{\documentclass} stehen in \autoref{latex-uebersicht}.

Als Nächstes werden die Tagesordnungspunkte in der Umgebung \texttt{Toplist} aufgelistet:

\begin{Codeblock}[title={Liste der Tagesordnungspunkte}]{latex}
	\begin{Toplist}
		\item Erster Tagesordnungspunkt
		\item Zweiter Tagesordnungspunkt
		\item ...
	\end{Toplist}
\end{Codeblock}

Danach wird jeder Tagesordnungspunkt in einer eigenen \texttt{Top}-Umgebung beschrieben:

\begin{Codeblock}[title={Erstellen eines Tagesordnungspunkts mit Abstimmungen}]{latex}
	\begin{Top}{Titel des ersten Tagesordnungspunkts}
		Hier die ausführliche Beschreibung und Diskussion zum ersten TOP.
		% Abstimmungen sind optional
		\Abstimmung[]{Erste Abstimmung zum ersten Tagesordnungspunkt}{2}{1}{2}
		\Abstimmung[vertagt]{Zweite Abstimmung zum ersten Tagesordnungspunkt}{}{}{}
	\end{Top}
\end{Codeblock}

Innerhalb eines Tagesordnungspunkts können eine oder mehrere Abstimmungen festgehalten werden. Dazu dient der Befehl \mintinline{latex}{\Abstimmung[Vermerk]{Gegenstand}{Ja}{Nein}{Enthaltungen}}. Bei einer regulären Abstimmung bleiben die eckigen Klammern leer und alle drei Stimmfelder werden mit einer Zahl gefüllt, auch wenn diese 0 ist.

Wurde eine Abstimmung nicht durchgeführt, etwa weil sie vertagt wurde, wird statt der Stimmen ein Vermerk eingetragen: Der Vermerk steht in den eckigen Klammern, die drei Stimmfelder bleiben leer (siehe die zweite Abstimmung im Beispiel). Wie das Ergebnis einer Abstimmung bestimmt wird, ist in \autoref{latex-fsrprotocol} beschrieben.

Der Gegenstand einer Abstimmung sollte möglichst im Wortlaut des Antrags festgehalten werden, da das Protokoll als Nachweis für die Beschlüsse dient.

Zum Schluss muss das Dokument durch Beenden der \texttt{document}-Umgebung abgeschlossen werden:

\begin{Codeblock}[title={Abschließen des Dokuments}]{latex}
	\end{document}
\end{Codeblock}

\subsubsection{Abschließen und Ablegen}
\label{protokoll-ablegen}
\begin{enumerate}
	\item Das fertige Protokoll \hyperref[latex-kompilieren]{kompilieren}.
	\item Die erzeugte PDF-Datei auf Vollständigkeit und Tippfehler prüfen und gegebenenfalls korrigieren und neu kompilieren.
	\item Die PDF-Datei aus dem Ordner \texttt{build/} in den Jahresordner verschieben, sodass \texttt{.tex}- und \texttt{.pdf}-Datei zusammen abgelegt sind.
\end{enumerate}

Bei Unklarheiten zur Erstellung eines Protokolls hilft der aktuelle Administrator weiter. Auch die \texttt{.tex}-Dateien älterer Protokolle können als Orientierung hilfreich sein.
